\subsection{Монотонные последовательности. Критерий сходимости монотонной последовательности.}

\textbf{Определения.}
\begin{enumerate}
	\item $\{a_n\} \uparrow$ — \textbf{строго возрастает}, если $\forall n\; a_n < a_{n + 1}$
	\item $\{a_n\} \uparrow$ — \textbf{возрастает (не убывает)}, если $\forall n\; a_n \leq a_{n + 1}$
	\item $\{a_n\}\downarrow$ — \textbf{строго убывает}, если $\forall n\; a_n > a_{n + 1}$ % ИСПРАВЛЕНО (< на >)
	\item $\{a_n\}\downarrow$ — \textbf{убывает (не возрастает)}, если $\forall n\; a_n \geq a_{n + 1}$
\end{enumerate}
\textbf{Определение.} $\{a_n\}$ — монотонна, если выполняется одно из условий 1-4.\\

\textbf{Теорема.} $\{a_n\} \uparrow$ и ограничена сверху $\Rightarrow a_n$ сходится и $a_n \to \sup a_n$\\
\textbf{Теорема.} $\{a_n\} \downarrow$ и ограничена снизу $\Rightarrow a_n$ сходится и $a_n \to \inf a_n$

\begin{tcolorbox}[title=Доказательство (случай конечного предела)]
	\begin{enumerate}
		\item $a_n \uparrow$ и ограничена сверху $\Rightarrow$ (по полноте $\R$) $\exists \sup a_n := A$
    \begin{enumerate}
        \item $\forall n\;\; a_n \leq A$
        \item $\forall \eps > 0\; \exists n_\eps:\; a_{n_\eps} > A - \eps$\\
            $\forall n > n_\eps\;\; A - \eps < a_{n_\eps} \leq a_n \leq A < A + \eps \Rightarrow |a_n - A| < \eps$, т.е. $a_n \to A$
    \end{enumerate}
		\item $a_n \downarrow$ и ограничена снизу $\Rightarrow$ (по полноте $\R$) $\exists \inf a_n := a$ % ИСПРАВЛЕНО (полносте -> полноте)
			\begin{enumerate}
				\item $\forall n\;\; a_n \geq a$
				\item $\forall \eps > 0\; \exists n_\eps:\; a_{n_\eps} < a + \eps$\\
					$\forall n > n_\eps\;\; a-\eps < a \leq a_n \leq a_{n_\eps} < a + \eps \Rightarrow |a_n - a| < \eps$, т.е. $a_n \to a$
			\end{enumerate}
	\end{enumerate}
\end{tcolorbox}

\textbf{Теорема.} Монотонная последовательность сходится $\Leftrightarrow$ она ограничена. \\
\textbf{Теорема (Обобщенная).} Монотонная последовательность всегда имеет предел (конечный или бесконечный).

\begin{tcolorbox}[title=Доказательство (общий случай)]
	\begin{enumerate}
		\item Если $\{a_n\}$ монотонна и ограничена $\Rightarrow$ (см. выше) $\{a_n\}$ имеет конечный предел.
		\item Пусть $\{a_n\}$ монотонна и неограничена.
			\begin{enumerate}
				\item Пусть $\{a_n\} \uparrow$. Так как $\{a_n\}$ не ограничена (и ограничена снизу первым членом), то она не ограничена сверху.\\
				$\forall E > 0 \; \exists n_E:\; a_{n_E} > E$. \\
				В силу возрастания: $\forall n > n_E \;\;\; a_n \geq a_{n_E} > E$. \\
				Это означает, что $a_n \to +\infty$.
				
				\item Пусть $\{a_n\} \downarrow$. Так как $\{a_n\}$ не ограничена, то она не ограничена снизу.\\
				$\forall E > 0 \; \exists n_E:\; a_{n_E} < -E$. \\
				В силу убывания: $\forall n > n_E \;\;\; a_n \leq a_{n_E} < -E$. \\
				Это означает, что $a_n \to -\infty$.
			\end{enumerate}
	\end{enumerate}
\end{tcolorbox}
